\chapter[Introduction]{Introduction}
\label{cp:introduction}

% Writing guidance: This chapter follows a funnel structure — broad context (supply
% chain risk is a growing threat), narrowing to the specific gap (existing SBOM tools
% treat inventories as flat lists, not graphs amenable to structural analysis), then
% zooming into what was built (SBOM Lens) and what was found. Target length: 7–10 pages.
% Technical terms (CycloneDX, CVSS, articulation points) should be named here but NOT
% explained — forward-reference to Chapter 2 for definitions.


\section{Motivation}
\label{sec:motivation}

% Key point 1: Modern software is built on third-party open-source dependencies.
%   The scale is significant — cite a statistic from the Synopsys OSSRA report or
%   similar (e.g., "X% of commercial codebases contain open-source components").
%   Frame the software supply chain as the aggregate of all these dependencies and
%   the trust relationships between them.

% Key point 2: Supply chain attacks exploit that trust. Three attacks establish the stakes:
%   - SolarWinds SUNBURST (December 2020, CISA Alert AA20-352A): attackers compromised
%     the build pipeline of a trusted vendor; malicious update propagated to ~18,000
%     organisations including US government agencies. Demonstrates cascading failure
%     through vendor trust relationships.
%   - log4shell (CVE-2021-44228, December 2021): a critical flaw in the Apache Log4j
%     library — a transitive dependency in millions of Java applications — exposed an
%     estimated 3 billion devices. Demonstrates how a single indirect dependency can
%     compromise an entire ecosystem.
%   - xz-utils backdoor (CVE-2024-3094, CVSS 10.0, March 2024): a two-year social
%     engineering campaign inserted a backdoor into a core Linux compression utility,
%     narrowly avoided widespread SSH compromise. Demonstrates patience and stealth
%     in supply chain attacks.

% Key point 3: Policy response — the SBOM mandate. US Executive Order 14028
%   (May 2021) mandated that software sold to the US federal government must be
%   accompanied by a Software Bill of Materials (SBOM) — a machine-readable inventory
%   of all components and dependencies. EO 14028 Section 10(j) provides the authoritative
%   definition. EO 14144 (January 2025) reinforced and extended these requirements.
%   NIST and CISA have published SBOM minimum elements guidance. Cite: \citep{EO14028}.

% Key point 4: The transparency gap. Despite SBOM adoption, existing analysis tools
%   (Dependency-Track, grype, syft, Sunshine/CycloneDX) treat SBOMs as flat component
%   inventories. They answer "what components are present?" and "do any have known CVEs?"
%   but do not model the dependency graph structure. They cannot identify which components
%   are structural bottlenecks, how vulnerabilities would propagate through the dependency
%   network, or how to prioritise remediation based on graph position.

% Key point 5: The opportunity — graph theory and network science. A dependency graph
%   is a directed graph by nature. Network science provides a rich toolkit for analysing
%   such graphs: articulation points reveal single points of failure; k-core decomposition
%   identifies the most deeply embedded components; spectral analysis and epidemic models
%   estimate vulnerability propagation potential. These techniques are mature in other
%   domains (biological networks, social networks, internet topology) but have not been
%   systematically applied to SBOM dependency graphs.

% Key point 6: Thesis statement (one paragraph, written directly — not a comment).
This dissertation presents SBOM Lens, a graph-based framework for software supply chain risk assessment. The system ingests Software Bills of Materials in CycloneDX format and constructs a directed dependency graph, over which network science techniques are applied to surface structural vulnerabilities not visible through flat-inventory analysis. The design, implementation, and evaluation of SBOM Lens are described in \autoref{cp:system-design}, \autoref{cp:vuln-analysis}, and \autoref{cp:network-sci}, respectively, with empirical results presented in \autoref{cp:evaluation}.


\section{Problem Statement}
\label{sec:problem-statement}

% Key point 1: Define the specific problem precisely. Existing SBOM tooling has two
%   limitations: (a) vulnerability matching is done component-by-component without
%   considering the graph position of a vulnerable component — a vulnerability in a
%   deeply embedded, highly connected node is treated identically to one in a leaf node;
%   (b) no structural risk metrics are computed — concepts like articulation points,
%   k-core depth, or community structure are absent from existing tools.

% Key point 2: Consequence of the gap. Without graph-based analysis, security teams
%   must manually triage hundreds of CVEs from a flat list, with no way to prioritise
%   based on structural impact. A CVE in a shared transitive dependency used by 80% of
%   the graph's components is not distinguishable from one in an isolated leaf component.

% Key point 3: The research gap statement. Cite the ACM landscape study (2025) or
%   similar to ground the claim that graph-based SBOM analysis is absent from the
%   current tool ecosystem. Note GUAC (Graph for Understanding Artifact Composition,
%   OpenSSF 2024) as the closest related work — GUAC models provenance graphs across
%   projects, but does not apply network science metrics for within-SBOM risk assessment.
%   Forward-reference to \autoref{cp:related-work} for detailed comparison.

% Key point 4: The proposed solution in one sentence. SBOM Lens fills this gap by
%   constructing a per-SBOM directed dependency graph and applying a suite of network
%   science algorithms to produce graph-level and node-level risk metrics. The research
%   question is not "can this be done?" (the system exists) but "what does it reveal?"
%   (the empirical evaluation in \autoref{cp:evaluation} answers this).


\section{Research Objectives}
\label{sec:objectives}

% Writing guidance: Objectives are framed as design-and-analysis goals, not hypothesis
% questions, because this is a tool/system thesis with a built artefact. Each objective
% maps to one implementation or evaluation chapter — this alignment allows the Conclusion
% to verify each was addressed.

% The following four objectives should be written as a numbered list (\begin{enumerate}):

% Objective 1 — Graph construction pipeline (maps to Chapter 4: System Design):
%   "To design and implement a directed dependency graph construction pipeline that
%   ingests CycloneDX JSON SBOMs and produces an annotated dependency graph amenable
%   to further analysis."

% Objective 2 — Vulnerability integration (maps to Chapter 5: Vulnerability Analysis):
%   "To integrate multi-source vulnerability data (NVD, OSV, and GitHub Security
%   Advisories) with full CVSS v3.1 base score computation and overlay the results
%   on the dependency graph."

% Objective 3 — Network science algorithms (maps to Chapter 6: Network Science):
%   "To apply network science algorithms — including articulation point detection,
%   k-core decomposition, Louvain community detection, HITS hub/authority scoring,
%   spectral radius computation, and an SIR-analogy epidemic propagation model —
%   to quantify structural risk in software dependency graphs."

% Objective 4 — Empirical evaluation (maps to Chapter 7: Evaluation):
%   "To evaluate the framework against a set of real-world open-source project SBOMs
%   and demonstrate its ability to surface structural risk not visible through
%   flat-inventory analysis."


\section{Contributions}
\label{sec:contributions}

% Writing guidance: Each contribution is a verifiable deliverable — a module, pipeline,
% or demonstrated capability. Use active verbs: "we design," "we implement," "we
% demonstrate." Contributions are written as a numbered list (\begin{enumerate}) with
% \textbf{} for the contribution name, followed by a dash and one or two sentences.

% Contribution 1 — SBOM Graph Visualization (Chapter 4: System Design):
%   \textbf{SBOM Graph Visualization} — An interactive directed dependency graph
%   constructed from CycloneDX JSON SBOMs and rendered with Cytoscape.js, supporting
%   topology exploration, node inspection, and interactive editing of the dependency
%   structure. This contribution provides the foundational data structure upon which
%   all analysis is performed.

% Contribution 2 — Vulnerability Analysis (Chapter 5):
%   \textbf{Multi-Source Vulnerability Analysis} — A vulnerability integration pipeline
%   that queries NVD, OSV, and GitHub Security Advisories, deduplicates findings across
%   sources, and computes full CVSS v3.1 base scores. Results are overlaid on the
%   dependency graph, visually distinguishing vulnerable nodes by severity band.

% Contribution 3 — Structural Risk Assessment (Chapter 6: Network Science, first half):
%   \textbf{Structural Risk Assessment} — A suite of algorithms that identify structural
%   vulnerabilities in dependency graphs: articulation point detection (single points of
%   failure), Shannon entropy of the degree distribution (dependency concentration),
%   diamond dependency detection (version conflict risk), and k-core decomposition
%   (depth of component entrenchment).

% Contribution 4 — Graph Analysis (Chapter 6: Network Science, second half):
%   \textbf{Graph Analysis} — Application of Louvain community detection (ecosystem
%   modularity), HITS algorithm (hub and authority scoring for dependency centrality),
%   spectral radius computation (epidemic threshold estimation), and an SIR-analogy
%   epidemic propagation model (vulnerability spread simulation).

% Contribution 5 — Empirical Evaluation (Chapter 7):
%   \textbf{Empirical Evaluation} — Results on a set of real-world open-source project
%   SBOMs, reporting graph-level and node-level metrics and demonstrating that the
%   framework surfaces risk not visible through flat-inventory approaches. (Author note:
%   update with number of SBOMs once dataset is selected — see \autoref{cp:evaluation}.)


\section{Scope and Limitations}
\label{sec:scope}

% Key point 1: SBOM format scope. SBOM Lens supports CycloneDX JSON only. SBOMs in
%   SPDX or SWID formats are not supported. This is an explicit design boundary, not
%   an oversight — CycloneDX JSON is the format with the richest dependency graph
%   representation in the OWASP ecosystem.

% Key point 2: Evaluation scope. The evaluation is based on graph metrics computed on
%   real-world open-source project SBOMs. The thesis does not include a formal usability
%   study or user experiment. Assessment is metric-based, not user-experience-based.

% Key point 3: Feature scope — out of scope items. The following features are explicitly
%   out of scope for this thesis version:
%   - SBOM comparison / delta analysis (comparing two SBOMs for a project over time)
%   - Topology Optimizer (automated dependency graph restructuring recommendations)
%   - SBOMs for compiled binaries or container images (source-level SBOMs only)
%   The rationale for each exclusion should be stated briefly (scope management, not
%   fundamental limitations of the approach).

% Key point 4: Claims scope. The thesis demonstrates that graph theory and network
%   science metrics expose structural risk not visible in flat inventories. It does not
%   claim to solve software supply chain security comprehensively, nor does it claim
%   the algorithms are novel — the novelty lies in their systematic application to
%   SBOM dependency graphs.


\section{Thesis Structure}
\label{sec:structure}

% Writing guidance: Write one or two sentences per chapter stating the chapter's
% ARGUMENT or PURPOSE, not just its title. Anti-pattern: "Chapter 2 presents background."
% Good pattern: "Chapter 2 establishes the theoretical foundations..."
% Use \autoref{} for all chapter cross-references.

% The structure overview should be written as flowing prose (not a list), with one
% short paragraph per chapter. The following content must appear:

% \autoref{cp:introduction} (this chapter): Establishes the motivation for graph-based
%   SBOM analysis through three supply chain attack cases, identifies the gap in existing
%   tooling, and states the research objectives and contributions.

% \autoref{cp:background}: Establishes the theoretical and technical foundations
%   required to understand SBOM Lens — covering the SBOM concept and CycloneDX format,
%   software supply chain security, graph theory fundamentals, network science concepts,
%   and vulnerability databases (NVD, OSV, CVSS v3.1).

% \autoref{cp:related-work}: Surveys existing SBOM analysis tools (Dependency-Track,
%   grype, syft, Sunshine/CycloneDX, GUAC) and graph-based security research, then
%   identifies the gap that SBOM Lens addresses.

% \autoref{cp:system-design}: Describes the architecture of SBOM Lens — the React +
%   FastAPI system design, the CycloneDX JSON parsing pipeline, directed graph
%   construction, and the interactive SBOM Graph visualisation built with Cytoscape.js.

% \autoref{cp:vuln-analysis}: Presents the multi-source vulnerability integration
%   pipeline (NVD, OSV, GitHub Security Advisories), the CVSS v3.1 base score computation
%   model, vulnerability management workflow, and graph overlay visualisation.

% \autoref{cp:network-sci}: Presents the structural risk assessment algorithms
%   (articulation points, entropy, diamond dependency detection, k-core decomposition)
%   and graph analysis methods (Louvain community detection, HITS, spectral radius,
%   SIR-analogy epidemic propagation), explaining the theoretical basis and supply chain
%   risk interpretation of each.

% \autoref{cp:evaluation}: Evaluates the framework on a set of real-world open-source
%   project SBOMs, reporting graph metrics, vulnerability findings, and network science
%   results. Discusses what the metrics reveal and the limitations of the evaluation.

% \autoref{cp:conclusion}: Summarises the contributions, maps them back to the research
%   objectives stated in this chapter, acknowledges limitations of the current work, and
%   proposes concrete directions for future research.
